\begin{ledgroupsized}[r]{120mm}
\footnotesize
\pstart
\noindent\textbf{Überlieferung:}
\pend
\end{ledgroupsized}
\begin{ledgroupsized}[r]{114mm}
\footnotesize
\pstart
\parindent -6mm
\makebox[6mm][l]{\textit{L}}%
Aufzeichnung:
LH~XXXV~10, 6~Bl.~38.
Ein Zettel (ca.~9,5~x~10,5~cm.);
Ränder nach Abschnitt unregelmäßig.
Zwei Seiten.
\pend
\end{ledgroupsized}
\selectlanguage{latin}
\vspace{8mm}
\pstart
\noindent
\href{\Scan{LH_35_10_06_038r}}%
{\lbrack38~r\textsuperscript{o}\rbrack}\ 
\edtext{Theorema I. Pondera\indexS{pondus regulam retinens} 
quae retinent regulas\indexS{regula inclinata}
sunt inter se ut secantes angulorum\indexS{secans anguli regulae ad perpendicularem} 
quos faciunt regulae ad perpendicularem.}{%
\lemma{}%
\Bfootnote{Theorema I. \lbrack...\rbrack\ perpendicularem. \textit{erg.~L}}}
\pend
\begin{edfigure}[0.5em][1em]
\hfill
\includegraphics[scale=1]{%
gesamttex/edit_VIII-4/images/LH_35_10_06_038_d1_038r.pdf}
\hfill%
\includegraphics[scale=1]{%
gesamttex/edit_VIII-4/images/LH_35_10_06_038_d2_038r.pdf}
\hfill
\edcaption{\hspace*{3mm}%
\lbrack\textit{Fig.~1}\rbrack 
\hspace*{48mm}%
\lbrack\textit{Fig.~2}\rbrack}
\end{edfigure}
\pstart
\edtext{\textls{Theorema} II. Si alia regu\textlangle la\textrangle\ sit perpendicularis\indexS{regula perpendicularis} pondera sunt aequalia.
Coincidit \textit{Q} et \textit{M}\lbrack,\rbrack\ \textit{P} et \textit{N}.}{%
\lemma{}%
\Bfootnote{\textls{Theorema} \lbrack...\rbrack\ et \textit{N}. \textit{erg.~L}}%
\lemma{\textls{Theorema} II. \lbrack...\rbrack\ et \textit{N}}%
\Cfootnote{Das Theorem samt \lbrack\textit{Fig.~3}\rbrack\ wird im weiteren Verlauf von N.\;\ref{60030} nicht mehr behandelt. Siehe aber die ausführliche Besprechung in N.\;\ref{60100} (als "Theorema 2", \refpassage{35_10_13_001-002_8}{35_10_13_001-002_7b}).}}
\pend
\begin{edfigure}[0.1em][1em]
\includegraphics[scale=1]{%
gesamttex/edit_VIII-4/images/LH_35_10_06_038_d3_038r.pdf}
\edcaption[0pt]{\lbrack\textit{Fig.~3}\rbrack}
\end{edfigure}
\pstart
\edtext{Globus gravis\indexS{globus gravis}\indexS{globus duas rectas rigidas tangens} duas}{%
\lemma{Globus}%
\Bfootnote{%
\bars{\textit{A} \textit{gestr.}} gravis
\textit{(1)}~tangit simul
\textit{(2)}~duas~\textit{L}}}
rectas\indexS{recta rigida}
\edtext{rigidas}{\lemma{}\Bfootnote{rigidas \textit{erg.~L}}}
\textit{BC}, \textit{DE}, in eodem plano inter 
\edtext{se ejus}{%
\lemma{se}%
\Bfootnote{%
\textit{(1)}~et cum
\textit{(2)}~ejus~\textit{L}}}
centro \textit{A} existentes simul tangit et premit in punctis \textit{F} 
\edtext{et \textit{G}. Quaeritur qua vi\indexS{vis urgendi rectae rigidae} urgeantur hae rectae. Seu,}{%
\lemma{et \textit{G}.}%
\Bfootnote{%
\textit{(1)}~Unde ducantur
\textit{(2)}~Quaeritur qua vi
\textit{(a)}~plana
\textit{(b)}~urgeantur hae rectae.
\textit{(aa)}~Seu~(si
\textit{(bb)}~Ponantur
\textit{(cc)}~Seu,~\textit{L}}}
si ponantur ita coercitae ut non cedere possint nisi in perpendicularibus ad semet, 
quanta debeant esse pondera\indexS{pondus cessionem impediens}
\textit{K} et \textit{I}, quae cessionem eorum 
\edtext{impediant. Jungantur}{%
\lemma{impediant.}%
\Bfootnote{%
\textit{(1)}~Ex puncto \textit{A} agatur \bars{ad \textit{BC} \textit{streicht Hrsg.}} ipsa \textit{AH} parallela ipsi \textit{DE}, et ex eodem \bars{puncto \textit{streicht Hrsg.}}
\textit{(a)}~\textit{A}
\textit{(b)}~agatur ad \textit{DE}, ipsa \textit{AL} parallela rectae \textit{BC}. Dico
\textit{(2)}~Jungantur
\textit{(3)}~Ex A jung
\textit{(4)}~Jungantur~\textit{L}}}
\textit{AF} et \textit{AG}, et agatur ad \textit{BC}, ipsa \textit{AH} et ad \textit{DE} ipsa \textit{AL}, ita ut \textit{AH} et \textit{AL} sint 
\edtext{parallelae illa}{%
\lemma{parallelae}%
\Bfootnote{%
\textit{(1)}~ipsis
\textit{(2)}~illa~\textit{L}}}
ipsi \textit{DE}, haec 
\edtext{ipsi \textit{BC}. Jungatur $HQ=AG$, et parallela, et $LP=AC$ et parallela et perpendicularis per \textit{A}, secet \textit{HQ} in \textit{M} et \textit{LP} in \textit{N}. Patet}{%
\lemma{}%
\Bfootnote{%
ipsi \textit{BC}.
\bars{%
\textit{(1)}~Denique recta perpendicularis per \textit{A}, secet \textit{DE} in \textit{N} et \textit{BC} in \textit{M}.
\textit{(2)}~Jungatur $HQ=AG$, \bars{et parallela \textit{erg.}}, et \lbrack...\rbrack\ in \textit{N}.~\textit{erg.}} 
Patet~\textit{L}}}
corpus \textit{A} agere in \textit{BC},
\edtext{quasi in recta inclinata\indexS{recta inclinata} \textit{AH} descendens,}{%
\lemma{quasi}%
\Bfootnote{%
\textit{(1)}~ferretur
\textit{(2)}~in recta inclinata
\textit{(a)}~\textit{DE}
\textit{(b)}~\textit{AH}
\textit{(aa)}~parallela ipsi \textit{D}
\textit{(bb)}~descendens,~\textit{L}}}
et nisum\indexS{nisus in rectam inclinatam} in \textit{AH} 
esse ad nisum absolutum\indexS{nisus absolutus} ut \textit{AH} 
\edtext{ad \textit{AM}, sed}{%
\lemma{ad \textit{AM},}%
\Bfootnote{%
\textit{(1)}~\bars{et \textit{streicht Hrsg.}}
\textit{(2)}~et
\textit{(3)}~sed~\textit{L}}}
actio\indexS{actio in rectam inclinatam} in \textit{AH} quae pervenit in \textit{BC}, est ut \textit{AF} ad \textit{AH}. Ergo pondus \textit{K} est ad 
\edtext{pondus \textit{A} datum,}{\lemma{}\Bfootnote{pondus \bars{\textit{A} \textit{erg.}} datum,~\textit{L}}}
ut \textit{AF} ad \textit{AM}, et pondus \textit{I} ad \textit{A} 
ut \textit{AG} 
\edlabel{35_10_06_038_1a}%
\edtext{}{%
\xxref{35_10_06_038_1a}{35_10_06_038_1b}%
\lemma{ad \textit{AN}.}%
\Bfootnote{%
\href{\Scan{LH_35_10_06_038v}}%
{\lbrack38~v\textsuperscript{o}\rbrack}\ 
\textit{(1)}~Sut\lbrack\textit{!}\rbrack\
\textit{(2)}~Sunt ergo
\textit{(3)}~\bars{Sed \textit{streicht Hrsg.}}
\textit{(a)}~\textit{AG} et \textit{A}
\textit{(b)}~\bars{\textit{AG} \textit{streicht Hrsg.}} et \textit{AF} sunt \bars{aequales radio \textit{streicht Hrsg.}}
\textit{(aa)}~ergo
\textit{(bb)}~et Rad. \textit{AG} ad rad. \textit{AF}
\textit{(aaa)}~sunt \bars{reciproce \textit{streicht Hrsg.}} ut \textit{AF} ad \textit{AG}. Ergo
\textit{(bbb)}~ergo reciproce sunt
\textit{(ccc)}~\bars{inter se \textit{streicht Hrsg.}} reciproce
\textit{(4)}~Nisus~\textit{L}}}%
ad \textit{AN}.
\pend
\pstart
\begingroup
\setlength\edtabcolsep{0.375em}
\begin{edtabularc}
\href{\Scan{LH_35_10_06_038v}}%
{\lbrack38~v\textsuperscript{o}\rbrack}\ 
&
Nisus\edlabel{35_10_06_038_1b} in
&
\textit{AF}
&
ad absolutum\indexS{nisus absolutus}
&
ut \textit{AF}
&
seu rad. globi
&
ad \textit{AM}
\\
&
\makebox[3.8em]{\dotfill}
&
\textit{AG}
&
\makebox[6.3em]{\dotfill}
&
ut \textit{AG}
&
\makebox[6.5em]{\dotfill}
&
ad \textit{AN}
\end{edtabularc}
\hfill
\endgroup
\pend
\newpage
\pstart
Nisus in \textit{AF} ad absolutum\indexS{nisus absolutus}
\edtext{ut rad.}{%
\lemma{ut}%
\Bfootnote{%
\textit{(1)}~\textit{AF} ad \textit{A}
\textit{(2)}~rad.~\textit{L}}}
ad \textit{AM}, absolutus\indexS{nisus absolutus} ad nisum in \textit{AG} ut \textit{AN} ad rad.
\pend
\pstart
Ergo nisus\indexS{nisus super regula}
\edtext{in \textit{AF} seu super regula \textit{BC} ad nisum\indexS{nisus super regula} in \textit{AG} seu super regula \textit{DE},}{%
\lemma{in \textit{AF}}%
\Bfootnote{%
\textit{(1)}~ad nis. in \textit{AG},
\textit{(2)}~seu \lbrack...\rbrack\ regula \textit{DE},~\textit{L}}}
ut 
$\dfrac{\text{rad.}}{AM}\;\dfrac{AN}{\text{rad.}}$%
\rule[-2.7ex]{0pt}{6.6ex}
seu ut \textit{AN} ad
\edtext{\textit{AM}, seu ut rectae}{%
\lemma{\textit{AM}, seu}%
\Bfootnote{%
\textit{(1)}~reciproce 
\textit{(1)}~ut rectae~\textit{L}}} 
occursus perpendicularis per centrum\indexS{centrum gravis} cum parallelis ad 
\edtext{radios ductis ex punctis occursuum parallelarum}{%
\lemma{radios}%
\Bfootnote{%
\textit{(1)}~ex
\textit{(2)}~ductis ex punctis
\textit{(a)}~parallelo
\textit{(b)}~occursuum parallelarum~\textit{L}}}
ad rectas rigidas\indexS{recta rigida} rectis rigidis occurrentium. 
\pend
\pstart
\edtext{Jam \lbrack posito radio eodem\rbrack\ \textit{AG}}{%
\lemma{Jam}%
\Bfootnote{%
\textit{(1)}~\textit{AN}
\textit{(2)}~\bars{posito radio eodem \textit{gestr., wieder gültig gemacht Hrsg.}} \textit{AG}~\textit{L}}}
vel \textit{AF}, patet \textit{AM} 
\edtext{esse ad \textit{AH}\lbrack,\rbrack\ \lbrack ut\rbrack\ secantem anguli\indexS{secans anguli regulae ad perpendiculum} \textit{HAM} ad radium, id est anguli}{%
\lemma{esse}%
\Bfootnote{%
\textit{(1)}~secantem anguli \textit{HAM}, id est anguli \bars{ad \bars{radium \textit{streicht Hrsg.}}~\textit{erg.}}
\textit{(2)}~ad \textit{AH} \bars{ut \textit{erg. Hrsg.}} secantem \lbrack...\rbrack\ anguli~\textit{L}}}
quem recta \textit{DE} facit ad perpendiculum, et similiter \textit{AN}
\edtext{esse ad \textit{AL},}{%
\lemma{esse}%
\Bfootnote{%
\textit{(1)}~secantem
\textit{(2)}~ad \bars{ad \textit{gestr.}} \textit{AL},~\textit{L}}}
ut secantem anguli\indexS{secans anguli ad radium} \textit{LAN} ad radium. 
Ergo
\begin{tabular}{c}
$AM=AH \cdot \sec{(HAM)}:\text{rad.}$
\\$AN=AL\, \lbrack\cdot\rbrack \sec{(LAN)}:\text{rad.}$
\end{tabular}
Ergo $AM:AN=$  in ratione composita \textit{AH} ad \textit{AL} et secantium, 
sed \textit{AH} est ad \textit{AF} ut 
\edtext{secans anguli\indexS{secans anguli ad radium}}{%
\lemma{secans}%
\Bfootnote{%
\textit{(1)}~complementi
\textit{(2)}~anguli~\textit{L}}}
\textit{HAF} ad radium, qui
\edtext{est angulus complementi\indexS{angulus complementi regularum}}{%
\lemma{}%
\Bfootnote{%
est \bars{qui angulus est\textit{str. Hrsg.}} angulus~\textit{L}}}
complementi ipsius 
\edtext{\lbrack\textit{AHF}\rbrack,}{%
\lemma{}%
\Bfootnote{%
\textit{FAH} \textit{L ändert Hrsg.}}}
quem faciunt regulae\indexS{angulus qui faciunt regulae inter se} inter
\edtext{se. Ergo}{%
\lemma{se.}%
\Bfootnote{%
\textit{(1)}~Et pari jure
\textit{(2)}~Ergo~\textit{L}}}
pari jure et
\edtext{\textit{AL} est secans\indexS{secans anguli complementi regularum} complementi regularum.}{%
\lemma{\textit{AL} est}%
\Bfootnote{%
\textit{(1)}~angulus com
\textit{(2)}~secans complementi
\textit{(a)}~planorum
\textit{(b)}~regularum.~\textit{L}}}
Nempe \textit{AL} est 
\edtext{ad \textit{AG} ut secans}{%
\lemma{ad \textit{AG}}%
\Bfootnote{%
\textit{(1)}~radium
\textit{(2)}~ut secans~\textit{L}}}
anguli \textit{LAG} ad rad. 
Sed ang. \textit{LAG} est compl. anguli \textit{ALG}, qui est regularum.\indexS{angulus regularum} 
Ergo \textit{AH} et \textit{AL} sunt aequales. 
Ergo pondera\indexS{pondus cessionem impediens} 
sunt ut secantes angulorum,\indexS{secans anguli regulae ad perpendiculum} 
quos regulae quae urgentur
\edtext{faciunt \lbrack ad\rbrack\ perpendiculum.}{%
\lemma{}%
\Bfootnote{%
faciunt \bars{ad \textit{gestr., wieder gültig gemacht Hrsg.}}
\textit{(1)}~pendicul
\textit{(2)}~perpendiculum.~\textit{L}}}
\pend
\pstart
\edtext{%
Non est necesse ut pondera \textit{I} et \textit{K} aequentur ipsi \textit{A}, sed
summa rectangulorum\indexS{summa ductorum ex celeritatibus in pondera} ex celeritatibus in 
\edtext{pondera sua}{%
\lemma{pondera}%
\Bfootnote{%
\textit{(1)}~aequatur
\textit{(2)}~sua~\textit{L}}}
\textit{I} et \textit{K} aequatur 
\edtext{rectangulo\indexS{ductus ex pondere in celeritatem} ex}{%
\lemma{rectangulo}%
\Bfootnote{%
\textit{(1)}~ex celeritate in pondus \textit{A}
\textit{(2)}~ex~\textit{L}}}
pondere \textit{A} in suam celeritatem.}{%
\lemma{}%
\Bfootnote{%
Non est \lbrack...\rbrack\ celeritatem. \textit{erg.~L}}}
\pend
